\documentclass[10pt,a4paper]{amsart}
\usepackage[margin=19mm]{geometry}
\usepackage{amsmath,amssymb,mathtools}
\usepackage[scr=rsfs,cal=euler]{mathalfa}
\usepackage{xfrac}
\usepackage[unicode,hidelinks,bookmarks=false]{hyperref}
\newcommand{\ie}{i.e.\ }
\newcommand{\cf}{cf.\ }
\newcommand{\resp}{respectively\ }
\newcommand{\Subsection}{\subsection}
\providecommand{\nobreakdash}{\nobreak\hbox{-}}
\setlength{\emergencystretch}{2em}
\multlinegap0pt
\usepackage{bm}
\usepackage{bbm}
\usepackage{dsfont}
\usepackage{xspace}
\usepackage{cancel}
\usepackage{mathtools}
\usepackage{amsthm}
\makeatletter
\@ifundefined{theorem}{\newtheorem{theorem}{Theorem}}{}
\@ifundefined{lemma}{\newtheorem{lemma}[theorem]{Lemma}}{}
\@ifundefined{proposition}{\newtheorem{proposition}[theorem]{Proposition}}{}
\makeatother
\title[Characterising Solutions of Anomalous Cancellation]{Characterising Solutions of Anomalous Cancellation}
\author{Satvik Saha, Sohom Gupta, Sayan Dutta, Sourin Chatterjee}
\date{Source: arXiv:2302.00479v1 (CC BY 4.0)}
\begin{document}
\maketitle

\noindent\emph{Source and licence.} Characterising Solutions of Anomalous Cancellation. Version arXiv:2302.00479v1 (CC BY 4.0), distributed under the Creative Commons Attribution 4.0 International licence, as recorded in the arXiv licence metadata for this exact version and shown on the abstract page https://arxiv.org/abs/2302.00479v1. Full rights notice: licenses/arxiv-2302.00479-CC-BY-4.0.txt.\par\smallskip

\noindent\textit{Editorial scope.} Counted unit: the Lemma whose source label is \texttt{lemma:gcds} in the article (source0.tex lines 755--760, source label \texttt{lemma:gcds}), delivered together with the article's own complete original proof of it (source0.tex lines 761--810, one \texttt{proof} environment of 50 lines). No mathematics is written, repaired or re-derived: statement and proof are the article's own text line for line. Required context retained uncounted: proposition:extension (lines 240--246); lemma:trivial\_2 (lines 492--501); lemma:last\_block (lines 669--676). Every definition, display and label of those lines is retained unchanged. Omitted from this package and not redistributed: 1--239, 247--491, 502--668, 677--754, 811--1043. Nothing inside a retained excerpt is omitted or altered: every retained body is delivered exactly as the article's lines are written, so no omission receipt of any kind accompanies this package. Citations: the retained excerpts carry no citation command at all, so no bibliography entry of the article is transcribed and no reference is redistributed. Reference adaptation: none was needed and none was made; every reference inside the delivered unit resolves inside it, so no unresolved reference or citation is produced. Printed slips kept literal: no printed word, symbol, hypothesis or number of the counted statement or of its proof was altered. Layout and class: the article's own class is \texttt{sn-jnl} with packages including cancel, hyperref; the wrapper is the lane's pinned \texttt{amsart} editorial wrapper at 10pt on A4, which restates the theorem-like environments the retained text needs and adds the article's own macro definitions that the retained text needs (none), with extra wrapper packages (bm, bbm, dsfont, xspace, cancel, mathtools). The delivered numbering is the unit's own. Assets: no retained span contains a figure environment, a graphics-inclusion command, a PDF-page inclusion, a table or a TikZ picture; no image file is copied into this package, the package declares no asset dependency and no image file is redistributed with it. Licence: version arXiv:2302.00479v1 of this article is distributed under the Creative Commons Attribution 4.0 International licence, as recorded in the arXiv licence metadata for this exact version and shown on the abstract page https://arxiv.org/abs/2302.00479v1. A full rights notice is delivered at licenses/arxiv-2302.00479-CC-BY-4.0.txt. Characters: every retained excerpt is pure ASCII, so the delivered engine, XeLaTeX with its default Latin Modern text font, reports no missing character.
\begin{proposition} \label{proposition:extension}
    Let $B$ be an arbitrary integer base, and let \[
        N = [a_1a_2\dots a_\ell \,b\, c_1c_2\dots c_k], \qquad
        N^+ = [a_1a_2\dots a_\ell \,b\: b\: b\, c_1c_2\dots c_k].
    \] Then, $N$ has property $P_{\ell; k}$ if and only if $N^+$ has property
    $P_{\ell + 1; k + 1}$. We say that $N^+$ is an \emph{extension} of $N$.
\end{proposition}

\begin{lemma} \label{lemma:trivial_2}
    Let $N_k$ have property $P_k$ and at least one of the following hold.
    \begin{enumerate}
        \item $a = c$.
        \item $a_i = b$ for all $1 \leq i \leq k$.
        \item $c_i = b$ for all $1 \leq i \leq k$.
    \end{enumerate}
    Then, all the digits $a_i = c_i = b$, i.e.\ $N$ is a trivial solution for 
    the $P_k$ problem.
\end{lemma}

\begin{lemma} \label{lemma:last_block}
    Let $B$ be an arbitrary integer base and let $N_k$ have property 
    $P^*_k$.
    Then, the digits in the last block satisfy $c_k < c_{i < k} = b$.
    In other words, all solutions of $P^*_k$ look like \[
        N = [a_1a_2\dots a_k \,b\, b b\dots b c_k]
    \]
\end{lemma}

\begin{lemma}\label{lemma:gcds}
    Let $B$ be an arbitrary integer base and let $N_k$ have property 
    $P^*_k$. Then, $\gcd(c_k, B) > 1$, i.e.\ the last digit $c_k$ must share some
    factor $p > 1$ with $B$.
    Furthermore, if $a_k \neq b$, then $gcd(a_k - b, B) > 1$.
\end{lemma}

\begin{proof}
    Suppose that $N_k$ has property $P^*_k$. Then $a, b, c > 0$, $c_k > 0$, and \[
        (aB + b)c = a(bB^k + c), \qquad
        (ac - abB^{k - 1})B = (a - b)c.
    \] This gives $B \mid (a - b)c$; writing \begin{align*}
        a' &= [a_1a_2\dots a_{k - 1}], & a &= a'B + a_k, \\
        c' &= [c_1c_2\dots c_{k - 1}], & c &= c'B + c_k,
    \end{align*} 
    we have $B \mid (a - b)(c'B + c_k)$, hence $B\mid (a - b)c_k = (a'B + a_k - b)c_k$,
    hence $B \mid (a_k - b)c_k$. Let \[
        B = p_1^{\alpha_1}p_2^{\alpha_2}\dots p_m^{\alpha_m}
    \] be the prime factorization of $B$, with each $\alpha_i \geq 1$.
    Then, for each prime factor $p_i$ of $B$, $p_i^{\alpha_i}\mid (a_k - b)c_k$.
    Let $\beta_i$ be the greatest integer such that $p_i^{\beta_i} \mid c_k$.
    Then, we must have $p_i^{\alpha_i - \beta_i} \mid a_k - b$; if not, the greatest
    power of $p_i$ dividing $(a_k - b)c_k$ would have been strictly less than $\beta_i + 
    (\alpha_i - \beta_i) = \alpha_i$, a contradiction.

    It is clear that we cannot have all $\beta_i \geq \alpha_i$; if so, we would have all
    $p_i^{\alpha_i} \mid c_k$, hence the product $p_1^{\alpha_1}\dots p_m^{\alpha_m} = B
    \mid c_k$, but $0 < c_k < B$, a contradiction. Thus, there must be some $\beta_j <
    \alpha_j$ corresponding to which $p^{\alpha_j - \beta_j} \mid a_k - b$, hence 
    $\gcd(a_k - b, B) \geq p_j > 1$. This proves the second part of our lemma.

    To prove the first part, we use induction on the block size $k$.
    Consider $k = 1$, where $a_k = a$, $c_k = c$, and suppose that all $\beta_i = 0$. 
    This forces all $p_i^{\alpha_i - \beta_i} = p_i^{\alpha_i} \mid a - b$, hence 
    their product $p_1^{\alpha_1}\dots p_m^{\alpha_m} = B \mid a - b$. 
    But $0 \leq \vert a - b\vert < B$, forcing $a - b = 0$.
    By Lemma~\ref{lemma:trivial_2}, this gives a trivial solution, a contradiction.
    Thus, there must be some $\beta_{j'} > 0$, hence $\gcd(c, B) \geq p_j^{\beta_j'} > 1$.

    Next, suppose that the statement holds for some $k \geq 1$, and let \[
        N = [a_1a_2\dots a_k a_{k + 1} \,b\, c_1c_2\dots c_{k + 1}] =
        [a\,b\,c]
    \] have property $P^*_{k + 1}$. Then, $a, b, c, > 0$, $c_{k + 1} > 0$,
    and $B \mid (a_{k + 1} - b)c_{k + 1}$. Again, if all $\beta_i = 0$,
    then we must have all $p_i^{\alpha_i - \beta_i} = p_i^{\alpha_i} \mid a_{k + 1} - b$,
    hence their product $B \mid a_{k + 1} - b$.
    Since $0 \leq \vert a_{k + 1} - b \vert < B$, we have $a_{k + 1} - b = 0$. But we also have
    $b = c_1 > c_{k + 1}$ by Lemma~\ref{lemma:last_block}.
    By Proposition~\ref{proposition:extension}, the number with the digits $a_{k + 1}, 
    c_1 = b$ removed, i.e.\ \[
        N' = [a_1a_2\dots a_k\, b\, c_2\dots c_{k + 1}],
    \] must have the $P^*_k$ property (each block remains non-zero, and $b > 
    c_{k + 1}$ so not all digits are equal). Applying our induction hypothesis, we have 
    $\gcd(c_{k + 1}, B) > 1$.

    Thus, by induction, our statement holds for all $k \geq 1$.
\end{proof}

\end{document}
